% Options for packages loaded elsewhere
\PassOptionsToPackage{unicode}{hyperref}
\PassOptionsToPackage{hyphens}{url}
\PassOptionsToPackage{dvipsnames,svgnames,x11names}{xcolor}
%
\documentclass[
  12pt,
]{article}

\usepackage{amsmath,amssymb}
\usepackage{iftex}
\ifPDFTeX
  \usepackage[T1]{fontenc}
  \usepackage[utf8]{inputenc}
  \usepackage{textcomp} % provide euro and other symbols
\else % if luatex or xetex
  \usepackage{unicode-math}
  \defaultfontfeatures{Scale=MatchLowercase}
  \defaultfontfeatures[\rmfamily]{Ligatures=TeX,Scale=1}
\fi
\usepackage{lmodern}
\ifPDFTeX\else  
    % xetex/luatex font selection
\fi
% Use upquote if available, for straight quotes in verbatim environments
\IfFileExists{upquote.sty}{\usepackage{upquote}}{}
\IfFileExists{microtype.sty}{% use microtype if available
  \usepackage[]{microtype}
  \UseMicrotypeSet[protrusion]{basicmath} % disable protrusion for tt fonts
}{}
\makeatletter
\@ifundefined{KOMAClassName}{% if non-KOMA class
  \IfFileExists{parskip.sty}{%
    \usepackage{parskip}
  }{% else
    \setlength{\parindent}{0pt}
    \setlength{\parskip}{6pt plus 2pt minus 1pt}}
}{% if KOMA class
  \KOMAoptions{parskip=half}}
\makeatother
\usepackage{xcolor}
\usepackage[margin=1in]{geometry}
\setlength{\emergencystretch}{3em} % prevent overfull lines
\setcounter{secnumdepth}{3}
% Make \paragraph and \subparagraph free-standing
\ifx\paragraph\undefined\else
  \let\oldparagraph\paragraph
  \renewcommand{\paragraph}[1]{\oldparagraph{#1}\mbox{}}
\fi
\ifx\subparagraph\undefined\else
  \let\oldsubparagraph\subparagraph
  \renewcommand{\subparagraph}[1]{\oldsubparagraph{#1}\mbox{}}
\fi


\providecommand{\tightlist}{%
  \setlength{\itemsep}{0pt}\setlength{\parskip}{0pt}}\usepackage{longtable,booktabs,array}
\usepackage{calc} % for calculating minipage widths
% Correct order of tables after \paragraph or \subparagraph
\usepackage{etoolbox}
\makeatletter
\patchcmd\longtable{\par}{\if@noskipsec\mbox{}\fi\par}{}{}
\makeatother
% Allow footnotes in longtable head/foot
\IfFileExists{footnotehyper.sty}{\usepackage{footnotehyper}}{\usepackage{footnote}}
\makesavenoteenv{longtable}
\usepackage{graphicx}
\makeatletter
\def\maxwidth{\ifdim\Gin@nat@width>\linewidth\linewidth\else\Gin@nat@width\fi}
\def\maxheight{\ifdim\Gin@nat@height>\textheight\textheight\else\Gin@nat@height\fi}
\makeatother
% Scale images if necessary, so that they will not overflow the page
% margins by default, and it is still possible to overwrite the defaults
% using explicit options in \includegraphics[width, height, ...]{}
\setkeys{Gin}{width=\maxwidth,height=\maxheight,keepaspectratio}
% Set default figure placement to htbp
\makeatletter
\def\fps@figure{htbp}
\makeatother

\usepackage{amsmath}
\usepackage{amssymb}
\usepackage{booktabs}
\usepackage{threeparttable}
\usepackage{threeparttablex}
\usepackage{caption}
\usepackage{subcaption}
\usepackage{float}
\usepackage{tablefootnote}
% --- One vertically-centered float per page (Tables and Figures section) ---
\makeatletter
\setlength{\@fptop}{0pt plus 1fil}
\setlength{\@fpbot}{0pt plus 1fil}
\makeatother
\renewcommand{\floatpagefraction}{0.01}
\renewcommand{\textfraction}{0.01}
\usepackage{tikz}
\usepackage{pdflscape}
\usepackage{etoolbox}
\usepackage{graphicx}
\usepackage{hyperref}
\usepackage{adjustbox}
\hypersetup{
  colorlinks  = true,
  linkcolor   = black,
  citecolor   = black,
  urlcolor    = black
}
\usepackage[authoryear, round]{natbib}
\preto\section{\clearpage}

\makeatletter
\def\maxwidth{\ifdim\Gin@nat@width>\linewidth\linewidth\else\Gin@nat@width\fi}
\def\maxheight{\ifdim\Gin@nat@height>\textheight\textheight\else\Gin@nat@height\fi}
\makeatother
\setkeys{Gin}{width=\maxwidth,height=\maxheight,keepaspectratio}
\makeatletter
\@ifpackageloaded{caption}{}{\usepackage{caption}}
\AtBeginDocument{%
\ifdefined\contentsname
  \renewcommand*\contentsname{Table of contents}
\else
  \newcommand\contentsname{Table of contents}
\fi
\ifdefined\listfigurename
  \renewcommand*\listfigurename{List of Figures}
\else
  \newcommand\listfigurename{List of Figures}
\fi
\ifdefined\listtablename
  \renewcommand*\listtablename{List of Tables}
\else
  \newcommand\listtablename{List of Tables}
\fi
\ifdefined\figurename
  \renewcommand*\figurename{Figure}
\else
  \newcommand\figurename{Figure}
\fi
\ifdefined\tablename
  \renewcommand*\tablename{Table}
\else
  \newcommand\tablename{Table}
\fi
}
\@ifpackageloaded{float}{}{\usepackage{float}}
\floatstyle{ruled}
\@ifundefined{c@chapter}{\newfloat{codelisting}{h}{lop}}{\newfloat{codelisting}{h}{lop}[chapter]}
\floatname{codelisting}{Listing}
\newcommand*\listoflistings{\listof{codelisting}{List of Listings}}
\makeatother
\makeatletter
\makeatother
\makeatletter
\@ifpackageloaded{caption}{}{\usepackage{caption}}
\@ifpackageloaded{subcaption}{}{\usepackage{subcaption}}
\makeatother
\ifLuaTeX
  \usepackage{selnolig}  % disable illegal ligatures
\fi
\usepackage[]{natbib}
\bibliographystyle{plainnat}
\usepackage{bookmark}

\IfFileExists{xurl.sty}{\usepackage{xurl}}{} % add URL line breaks if available
\urlstyle{same} % disable monospaced font for URLs
\hypersetup{
  colorlinks=true,
  linkcolor={blue},
  filecolor={Maroon},
  citecolor={Blue},
  urlcolor={Blue},
  pdfcreator={LaTeX via pandoc}}

\author{}
\date{}

\begin{document}

% Switch footnotes to symbols for title page only
\renewcommand{\thefootnote}{\fnsymbol{footnote}}

\thispagestyle{empty}

\begin{center}
\vspace*{3cm}

{\Large Internalizing Environmental Risk: Insurance Design and Firm Behavior in Hazardous Industries}

\vspace{1cm}

{\normalsize \href{https://www.kalebkjavier.com}{Kaleb K. Javier}\footnote{University of California, Berkeley, Department of Agricultural and Resource Economics. Email: \texttt{kalebkja@berkeley.edu}. Website: \texttt{kalebkjavier.com}. I thank Meredith Fowlie, Billy Pizer, James Sallee, Nano Barahona, Ben Handel, Matt Backus, and Cailin Slattery for their invaluable comments and feedback. I also thank the participants and discussants at Camp Resources XXXI and the AERE Summer Conference 2025 for helpful discussions. All remaining errors are my own.}}

\vspace{0.2cm}

{\normalsize [Month] 2026}

\end{center}

\vspace{.2cm}

\begin{center}
{\small \textbf{Abstract}}
\end{center}

\begin{quote}
{\small \noindent Public trust funds in thirty-eight U.S. states insure underground storage tank
(UST) leaks at a uniform flat fee, even though leak risk varies widely with observable
traits such as tank age and wall construction. This paper quantifies how replacing
flat-fee pooling with actuarially priced coverage changes firm behavior and environmental
outcomes. I exploit Texas's 1999 closure of its trust fund, which forced all owners into
a private market that sets premiums according to facility risk. A facility-year panel
covering 1990--2021 links EPA inventories with administrative microdata for Texas and
eighteen control states that retained public funds. I construct a make-model comparison
sample that matches Texas facilities to control-state counterparts on the observable risk
characteristics, tank age bin and wall construction type, that private insurers price,
ensuring the comparison group is balanced on the dimensions along which the treatment
effect is theoretically heterogeneous. Difference-in-differences and event-study estimates
show that risk-rated premiums increased annual closure rates by two to three percentage
points, shifted closures from market exit to on-site replacement with safer equipment,
and reduced the probability of a reported leak by about one percentage point, nearly the
entire pre-policy baseline for the high-risk single-wall cohort. The treatment effect is
concentrated among older and single-walled facilities, consistent with the
premium-gradient mechanism, and is near zero for young double-walled tanks. The findings
provide the first causal micro-evidence that cost-based environmental liability insurance
can internalize pollution risk more effectively than uniform public pooling, offering
clear guidance for states that are reconsidering the structure of UST insurance programs.}
\end{quote}

% Switch back to arabic numerals for the rest of the document
\renewcommand{\thefootnote}{\arabic{footnote}}
\setcounter{footnote}{0}

\newpage

\section{Introduction}\label{introduction}

Compulsory liability insurance is the primary instrument through which
environmental law makes hazardous firms internalize the risk of harm
they cause. Whether it actually does so depends entirely on how the
insurance is priced. Strict liability makes polluters legally
responsible for harm, but it fails when firms are judgment-proof, when
victims are diffuse, or when causation is unobservable. Financial
responsibility mandates --- requiring firms to carry insurance or post
bonds before they may operate --- are the statutory response to these
failures {[}\citet{shavell1986}; \citet{shavell2004};
\citet{shavell2007}; OPA 1990; CERCLA 1980; RCRA 1976{]}. When firms
satisfy that obligation through compulsory insurance, the premium
becomes the operative price of environmental risk. A premium that rises
with the probability of harm leads the firm to internalize that risk
through ongoing costs. A premium invariant to risk satisfies the mandate
on paper but delivers no deterrence. Uniform per-tank premiums ensure
that hazardous firms carry coverage --- that is the good. Coverage
without risk pricing means the firm bears no marginal financial cost
from operating increasingly dangerous equipment --- that is the
distortion. The design of the premium, not the existence of the mandate,
determines whether compulsory insurance functions as an environmental
policy instrument. This paper provides the first firm-level causal
evidence on whether shifting compulsory insurance from a uniform
per-tank premium to actuarially differentiated, risk-rated pricing
changes what operators do with aging hazardous capital.

A uniform per-tank premium sets the marginal cost of operating an
additional aging, high-risk tank to zero. That removes the financial
incentive to retire it --- a textbook moral hazard distortion that
theory predicts and this paper confirms. The mechanism is precise. Under
a uniform per-tank premium \(\bar{P}\), the annual insurance cost is
\(\bar{P} \times N\), where \(N\) is the number of tanks. The marginal
cost of holding one more old, leak-prone tank is \(\bar{P}\) ---
identical to the cost of holding one new, low-risk tank. The operator
faces no financial gradient across the risk distribution of her
portfolio: \(\partial C / \partial \text{age} = 0\). Because the premium
is invariant to tank age and wall construction, she does not internalize
the age-hazard gradient. The efficient action --- retiring old
single-walled tanks before they leak --- is privately costless to avoid.
This is the hazard suppression this paper identifies and quantifies.
Uniform pricing also generates an adverse selection channel: it
cross-subsidizes high-risk facilities at the expense of low-risk ones,
crowds risk-rated private alternatives out of the market, and
concentrates the riskiest operators in pooled coverage --- the
environmental analog of the community-rating problem
\citep{pauly1974, rothschild1976}.
\citep{shavell1986, shavell2004, shavell2007} establish that compulsory
insurance restores deterrence if and only if premiums price observable
risk-relevant characteristics. For underground petroleum storage tanks,
tank age and wall construction type are exactly those characteristics
--- both verifiable by an underwriter at policy inception, and both
strong predictors of release probability. Whether risk-rated pricing in
practice disciplines high-risk operators through deterrence, drives them
out through adverse selection, or some combination, has not been
identified at the facility level. Separating these margins and
quantifying their welfare consequences requires a setting where the
transition from uniform to risk-rated pricing can be observed for a
large, well-measured population of facilities. Texas provides that
setting.

Underground petroleum storage tanks are the dominant source of
groundwater contamination in the United States. The harm is costly, and
the regulatory system --- not private tort law --- is the only operative
instrument for controlling it. Approximately 140,000 active retail
facilities and 550,000 active tanks operate nationwide \citep{epa2024}.
Roughly one in four facilities leaks over a 30-year operational life.
There have been approximately 550,000 confirmed releases since 1988.
Median cleanup cost is \$136,249 and mean cleanup cost is \$403,319 per
site (author's trust-fund claims data, \(N = 12{,}541\) sites across six
states); the right tail dominates expected loss. Beyond cleanup, benzene
and petroleum migrate silently into groundwater, reach drinking-water
wells, and generate health damages of approximately \$17,000 per release
\citep{marcus2021}. These harms generate no insurance claim and
therefore no actuarial signal. Private tort law does not reach this
problem. Contamination is slow, silent, and diffuse; by the time a plume
is detected it may have traveled hundreds of feet. Causation is
difficult to establish, and victims are dispersed and often unaware of
their exposure. The regulatory system is the only instrument. There is a
further structural feature with direct welfare implications. Actuarial
premiums price what insurers can recover --- cleanup costs that flow
through claims. They cannot price the health externality \(E\) because
it generates no third-party claim. Even a perfectly designed risk-rated
premium leaves \(E\) uninternalized. This is why a welfare gap between
risk-based pricing and the first-best Pigouvian instrument persists even
after reform. The combination of predictable risk, compulsory coverage,
externalized harm, and a pricing design that is a state-level policy
choice makes USTs the ideal laboratory for studying compulsory insurance
design.

Federal law mandates that every UST operator carry financial
responsibility coverage. But the mandate is silent on pricing design.
Thirty-eight states satisfied it through uniform per-tank premiums that
price no risk and eliminate the deterrence incentive the mandate was
intended to create. RCRA Subtitle I (1988) requires all UST operators to
demonstrate \$1 million in minimum financial responsibility coverage.
Compliance is enforced through operating permits. Federal law leaves the
pricing structure to states. States have chosen one of three approaches:
no state fund and reliance on the private market; a uniform per-tank
state trust fund; or risk-rated private insurance. Thirty-eight states
chose the trust fund model. In that model, every covered facility pays
the same annual per-tank assessment regardless of tank age, wall type,
or leak history. Operators are in full regulatory compliance, but
compliance is entirely decoupled from risk reduction. The policy
relevance extends well beyond Texas. More than a dozen of the 38 fund
states carry statutory sunset dates clustering through this decade:
Minnesota (2022), Colorado (2023), Kansas and Nebraska (2024), Illinois,
Missouri, and New Hampshire (2025), California and South Carolina
(2026). As these funds close, their facilities will face exactly the
transition Texas underwent in 1998. This paper measures the welfare
consequences of that transition. Texas provides the clean
identification: its uniform per-tank fund closed overnight on December
22, 1998, for reasons unrelated to the riskiness of any individual
facility, moving the entire state into risk-rated private insurance in a
single day.

On December 22, 1998, the Texas Petroleum Storage Tank Fund closed due
to aggregate fiscal insolvency. Approximately 26,000 active facilities
entered risk-rated private insurance overnight. The shock was externally
forced and institution-driven; its exogeneity comes from the fund's
balance sheet, not from any change in Texas facility risk profiles. The
insolvency was driven by aggregate claim volume across the entire fund,
not by the riskiness of any individual facility. Tank age, wall
construction, and installation cohort --- the characteristics that drive
release probability --- were unchanged at the reform date. The price of
risk changed; the underlying risk did not. Before December 22, 1998,
Texas facilities faced a uniform per-tank premium with
\(\partial P / \partial \text{age} = 0\). After, premiums are
actuarially differentiated, rising with tank age, wall construction
type, and piping technology in a schedule that closely tracks the
actuarial risk gradient (see \autoref{sec:data}). The control group
consists of 17 states that retained uniform per-tank coverage throughout
the study window and have complete administrative records on closure
dates, installation dates, and wall construction --- the exact
characteristics the private market prices. The panel spans 9.8 million
tank-years from 1990 to 2021 across 18 study states. I estimate the
causal effect using a difference-in-differences design with facility
fixed effects and cell-by-year fixed effects. Cells are defined by wall
type, fuel type, capacity, and installation cohort --- the dimensions
the private market prices and the dimensions along which the treatment
effect should be heterogeneous. The design delivers a within-cell,
within-facility, before-and-after comparison of Texas facilities against
observationally identical control-state facilities that never left
uniform per-tank coverage.

Risk-based pricing raises the annual tank closure rate by
\(\hat{\beta} = +1.58\) percentage points (s.e. = 0.40) on a pre-reform
control mean of 2.0\% --- a 79\% relative increase. The effect
concentrates precisely where pricing theory predicts: among old,
single-walled tanks where the risk-rated premium gradient is steepest. A
penalized logistic regression on 3.93 million pre-treatment
facility-years confirms the precondition. The annual first-release
hazard for single-walled tanks rises from 5.3 per 1,000 facility-years
in the youngest age bin to 12.3 per 1,000 in the oldest. For
double-walled tanks it remains flat. Observable characteristics carry
strong actuarial content, validating the premium gradient. The vintage
heterogeneity result provides the primary falsification test.
Pre-1989-vintage single-walled tanks show an ATT of \(+3.1\) pp;
post-1988 tanks show \(+0.6\) pp --- a fivefold differential. A
coincident Texas-wide economic shock would predict uniform closure
increases across all tank types. The theory predicts concentration among
old single-walled tanks where the premium gradient is steepest. The data
match the theory, not the alternative. Event-study pre-trends are flat,
with \(k = -1\) as the omitted reference period. The break is sharp at
the reform date and the elevated closure rate is sustained through 22
years post-reform. Annual LUST discovery falls \(-1.27\) pp (s.e. =
0.32) on a 1.73\% base, with both routine-monitoring and
closure-inspection channels declining. Tank retirement, not strategic
non-reporting, drives the environmental improvement. Among closing
facilities, the replacement share rises \(+6.3\) pp (s.e. = 1.3) and the
permanent-exit share falls \(-6.2\) pp (s.e. = 1.2). Risk-based pricing
induces fleet upgrading, not only market contraction. Translating these
results into welfare-comparable magnitudes and evaluating counterfactual
policies requires the structural model.

The difference-in-differences design identifies the average causal
effect of the 1998 reform. Three welfare-relevant questions lie beyond
its reach. First, the DiD delivers a closure-rate gap in percentage
points. Comparing alternative policies requires present-value welfare in
dollars, a translation the reduced form cannot make. Second, a
replacement subsidy and an age mandate were never implemented in Texas.
Their behavioral and welfare consequences are unobserved in any
reduced-form design and require model-based simulation. Third, actuarial
premiums price the cleanup costs that flow through insurance claims. The
health externality \(E \approx \$50{,}000\) per release generates no
insurance claim and is therefore structurally unpriced even under
risk-rated coverage. Quantifying how much of the first-best welfare gain
risk-based pricing fails to capture requires the structural primitives
--- specifically the estimated price-sensitivity parameter
\(\hat{\gamma}_p = -2.163\) (s.e. = 0.205) and the risk-internalization
parameter \(\hat{\gamma}_r = +0.114\) (s.e. = 0.005). That
\(\hat{\gamma}_r \ll \hat{\gamma}_p\) is itself the quantitative
statement of the residual wedge. These parameters imply the regime
ordering \(a^*_{\text{SOC}} < a^*_{\text{RB}} < a^*_{\text{UP}}\), where
\(a^*\) is the optimal tank-removal age threshold under each regime.
Risk-based pricing moves behavior toward the social optimum; the
unpriced externality prevents it from closing the gap. I estimate a
dynamic discrete-choice model of facility-level tank portfolio
management across a state space of 298,448 portfolios by Nested
Pseudo-Likelihood \citep{aguirregabiria2007}. I use the recovered
parameters to re-solve firm behavior and compute welfare under four
counterfactual policy instruments.

This paper contributes to three literatures: the theory and empirics of
environmental liability and financial responsibility; the empirical
literature on hazardous capital stock management under financial
assurance requirements; and the industrial organization of insurance
markets and welfare decomposition. On liability and financial
responsibility, \citep{shavell1986, shavell2004, shavell2007} provide
the theoretical foundation for compulsory insurance as a deterrence
instrument. \citet{yin2011a} and \citeauthor{yin2011b}
\citetext{\citeyear{yin2011b}; \citealp{yin2007}} study UST fund design
theoretically and with Michigan aggregate data. This paper provides the
first facility-level causal identification of how insurance pricing
design --- uniform versus risk-rated --- changes environmental behavior,
using a clean policy shock, within-facility panel identification, and a
structural welfare decomposition. On hazardous capital stock management,
\citet{boomhower2019} finds that bonding requirements change well
retirement and industry composition in oil and gas. \citet{armitage2026}
documents judgment-proof dynamics in aging well transfers. This paper
differs because the behavioral margin here is almost entirely extensive.
Wall-type retrofitting is prohibitively costly under 40 CFR Part 280
grandfathering rules --- a structural feature confirmed in the data:
only 0.003\% of active Texas facilities ever retrofitted tank
construction prior to terminal closure. The clean extensive margin
sharpens identification and simplifies the structural model. On the IO
of insurance markets, \citeauthor{einav2010}
\citetext{\citeyear{einav2010}; \citealp{einav2021}} and
\citet{bundorf2012} develop the welfare-decomposition logic for markets
with selection. A uniform per-tank premium is the environmental analog
of community rating in health insurance: premiums invariant to
observable risk create cross-subsidies and suppress efficient sorting.
But the uninternalized wedge here is not adverse selection on coverage.
It is the health externality that no actuarial premium can price. The
welfare decomposition reflects that distinction. On UST damages,
\citet{marcus2021} estimates groundwater health damages from UST leaks.
\citeauthor{guignet2014}
\citetext{\citeyear{guignet2014}; \citealp{guignet2018}},
\citet{walsh2017}, and \citet{zabel2012} estimate hedonic property-value
impacts. \citet{ho2018} studies UST remediation. This paper is
complementary: it studies the design of the insurance instrument meant
to prevent the damages those papers measure. The contribution is the
combination --- credible causal identification of behavioral responses
and a structural welfare decomposition that places those responses on a
policy-comparable scale --- that neither the liability literature nor
the hazardous capital literature has previously delivered in this
setting.

The paper proceeds as follows. \autoref{sec:setting} describes the
regulatory setting and institutional background, providing the
qualitative defense of identification. \autoref{sec:data} details data
construction from three administrative sources, sample construction, and
key descriptive facts. \autoref{sec:reduced-form} presents the
reduced-form evidence: the hazard model, the difference-in-differences
and event-study estimates on closure, leaks, and composition, and the
vintage heterogeneity analysis as the primary falsification test.
\autoref{sec:structural} develops the structural model: the portfolio
state space, nested logit payoffs, NPL estimation, and parameter
identification. \autoref{sec:welfare} uses the estimated model to
evaluate four counterfactual policies and reports the welfare
decomposition. \autoref{sec:conclusion} concludes.

\section{Setting}\label{setting}

An underground storage tank holds motor fuel at a gas station or a
commercial site. It sits in the ground for decades. A single-walled tank
is a bare shell. A double-walled tank adds a second wall and a monitored
space between the two. Single-walled tanks are cheaper to put in and
they leak more. The wall is set when the tank goes in the ground. To
change it, the operator has to dig the tank out and install a new one,
which costs about as much as starting fresh. In the Texas data, only
0.003\% of active facilities ever changed a tank's wall before closing
it. The wall is a fixed trait, not a dial an operator turns.

An operator who faces a higher price for risk has few ways to respond.
The cheap risk-reduction steps were already done. Federal deadlines in
1993 and 1998 required leak detection, spill protection, and corrosion
control, and almost every tank met them by the reform. A partial wall
upgrade does not help. Under 40 CFR Part 280, any such change
reclassifies the whole facility to current standards, and that cost
dwarfs any premium saving.\footnote{This is vintage-differentiated
  regulation (Gruenspecht 1982, Stavins 2005). Exempting existing
  capital from new standards extends the life of the oldest and most
  hazardous stock, because any modification triggers the current
  standards and that cost dwarfs the gain. It is the same mechanism as
  New Source Review in the Clean Air Act, where the threat of
  reclassification deterred investment at grandfathered plants (Bushnell
  2012, Heutel 2011). When new-capital rules cannot reach the existing
  stock, retirement-focused instruments are the right complement (Sallee
  2023, Jacobsen 2023).} Replacing a tank with a new double-walled one
works but is expensive, about \$55,000 a tank, and rare. That leaves one
lever that actually cuts risk: pull the old tank and close it. The
hazard is already in the ground, and it comes out only when the old
tanks do. So the margin I study is closure, and the policy question is
how to price keeping an old tank in service.

A typical site is one gas station with a handful of tanks. The country
has about {[}140,000{]} active facilities and {[}550,000{]} active
tanks. Sites differ in the fuel they sell and in how many tanks they
run. I use both later to see who responds to the reform.

The risk a tank carries is not hidden. It rises with two things anyone
can read off the tank: its age and its wall. Old tanks leak more than
new ones. Single-walled tanks leak more than double-walled ones. An
insurer can check both when it writes the policy. So it can charge a
riskier tank more. This is the fact the paper turns on, and I show it in
§4.1.

A leak is hard to catch and hard to pin on anyone. Fuel seeps into soil
and groundwater. It can travel for years before a drinking-water well or
a basement reveals it. By then the source is often unclear, and the
people harmed are spread out and may not know. Private lawsuits cannot
do much here. So the job of preventing leaks falls to regulation, not to
the courts.

A leak creates two kinds of cost. The first is cleanup, digging out and
treating the contaminated soil and water. Cleanup is expensive, with a
median around {[}\$136k{]} and a mean around {[}\$403k{]} per site. The
second is harm to people nearby. \citet{marcus_going_2021} shows that
UST leaks damage infant health, lowering birth weight among births near
a leak. Insurance pays for cleanup, because cleanup files a claim. It
does not pay for the health harm, because that harm files no claim. So
even a fairly priced premium leaves the health cost uncovered. That gap
is the wedge I return to in §6.

Federal law requires every tank operator to carry financial
responsibility for cleanup. The rule is RCRA Subtitle I, in force since
{[}1988{]}, with a minimum of {[}\$1 million{]} in coverage. Federal law
does not say how to price it. States chose one of three paths. Some let
operators buy risk-rated private insurance. Some set up a state trust
fund that charges every tank the same flat fee. Some did neither.
Thirty-eight states ran a flat-fee fund. A flat fee meets the mandate
but prices no risk. An old single-walled tank pays the same as a new
double-walled one.

Texas ran a flat-fee fund until it ran out of money. On December 22,
1998 the state closed the fund because it was insolvent, and about
{[}26,000{]} facilities moved to risk-rated private insurance at once.
The fund failed for a budget reason, not because any one facility got
riskier. So the price of risk changed for reasons outside any operator's
control. That is the shock I use. Texas went from a flat fee to a
risk-rated premium overnight, while the flat-fee control states did not.

\section{Data}\label{data}

I construct a facility-year panel of the underground storage tank
industry in Texas and 17 flat-fee control states, 1990--2021, by merging
four sources: a harmonized national tank registry, state leak-incident
records, Texas insurance-pricing data, and location and market measures
from the 2000 Census and facility geography. Appendix {[}X{]} documents
the construction.

The backbone of the panel is the EPA's LUST Finder, which harmonizes
state UST registries into one national schema
\citep{EPA2017USTCompendium}. My unit of observation is the
facility-year. A facility is a retail or commercial site with a unique
state registry number that operates one or more underground tanks. I
follow each facility from its first tank installation to its last
closure or to the end of 2021. The panel covers 117,250 facilities
across 18 states.

The registry has gaps in the fields a causal design needs. Of the 38
states with a uniform-premium trust fund, I keep the ones with complete
closure dates, install dates, and wall type. That drops 20 states and
leaves Texas and 17 controls. My sample is the tanks and facilities that
were open on the reform date, December 22, 1998. I follow each one
forward, year by year. A tank that closes later counts as a closure, not
a dropped observation. So closure and exit are outcomes I measure, not
cases the sample loses. I match Texas units to control units so I
compare like with like. The match is coarsened exact matching. At the
tank level I match on make-model and install year. At the facility level
I match on install year, facility size, and single-wall share. The
cell-by-year fixed effects in §4 control for these traits again.

Each registry record is a physical tank. It reports the tank's install
date, closure date, wall type (single-walled or double-walled), stored
fuel, and capacity. These attributes set the leak hazard and the
insurance premium. They also define the make-model cells used in §4.
From them I build each facility's stock variables: active tank count,
total capacity, average tank age, and single-walled share. I define two
fuel cuts for later. A gas station is any facility with a gasoline tank.
A motor-fuel-only facility is one whose tanks are all motor fuel. I also
sort facilities into four size groups by total capacity, G1 to G4.

My primary outcome is tank closure. It equals one in the year a tank is
permanently removed. At the facility I track the whole portfolio. A
facility can downsize, which means it closes some tanks but stays open.
It can also cut capacity, replace tanks, or exit. I split each closure
into two kinds. A replacement is a closure where the facility puts in a
new tank. A permanent exit is a closure where it does not. I use the
facility's full install history to tell them apart, so I do not count a
replacement as an exit.

Leak discovery comes from the state LUST incident files. Each record
gives a confirmed release by facility and report date. I split
discoveries by how they are found. Some show up in routine monitoring.
Others show up when a tank is closed and inspected. I drop states
without usable release dates when I pick the study states. So I measure
leak discovery for every state in the panel and do not need a separate
leak sample. One note on the dates. The 1995 move of state records into
the LUST Finder bunched install and closure dates near 1995. That
bunching is an administrative artifact, not real activity. I treat those
dates as artifacts, and I set the event-study reference year
(\(rel\_year = -1\)) after it.

For the heterogeneity analysis I attach measures fixed at the reform
that do not move with the policy. From the 2000 Census I classify each
facility as rural, low-population, low-income, or high-poverty. From
facility geography I build a market-thinness measure for gas retailers,
the count of active competitors within one mile in 1998. Because all of
these are set before or at the reform, they do not pick up the reform's
own effect on the market.

The insurance data let me measure the price each facility pays for risk.
For Texas I collected the filed rate manuals of the largest UST insurer
in the state through public-record requests to the Texas Department of
Insurance. I link these manuals to facilities through the Texas
Commission on Environmental Quality (TCEQ) database. That database
records each facility's insurance contract for about the past twenty
years. Texas is unusual in keeping this detail. It is what lets me
rebuild an annual premium for each Texas facility. For the control
states I use the ASTSWMO survey of state-fund fees. The survey gives the
flat per-tank fee each fund charged. The premium data have one main
limit. The Texas manuals cover the largest carrier, so the reconstructed
premium is the price under that carrier's filed schedule, not the
realized price across all carriers. I also collected trust-fund claims
data from 6 of the 17 control states. These records link LUST sites to
their cleanup claims. They are the basis for the cleanup-cost
distribution I show below.

\autoref{tbl-summary} reports summary statistics for the study
facilities at the reform date, with Texas and the control states in
separate columns. Panel A describes the capital stock. Texas facilities
are larger on average. They hold more three- and four-tank sites, while
the control states hold more single-tank sites
(\autoref{fig-baseline-tanks}). Texas capacity is also shifted toward
the larger bins (\autoref{fig-baseline-capacity}). The average Texas
facility holds {[}N{]} tanks and {[}N{]} gallons of capacity, against
{[}N{]} tanks and {[}N{]} gallons in the controls. Panel B reports
pre-reform annual rates. The closure rate is {[}N{]} percent and the
leak-discovery rate is {[}N{]} percent, and the two groups track each
other before the reform. \autoref{fig-cost-distribution} shows the
distribution of cleanup costs across LUST sites in the six claims
states. The median cleanup cost is {[}\$N{]} and the mean is {[}\$N{]},
which reflects a long right tail.

\section{Reduced-Form Evidence}\label{reduced-form-evidence}

This section presents the reduced-form evidence in five steps. I first
show that the traits an insurer can observe predict releases (§4.1). I
then show how the Texas market prices those traits, and how the Texas
premium pulled away from the flat control fee at the reform (§4.2). I
then state the design (§4.3). Section 4.4 asks whether the policy
worked: the average effect on tank closure, its timing, and whether
leaks fell. Section 4.5 asks who it worked on, where, and what
facilities did in response.

\subsection{Observable risk and the age-hazard
gradient}\label{observable-risk-and-the-age-hazard-gradient}

Risk-based pricing can only change behavior if the traits an insurer
sees at underwriting actually predict releases. They do, and the
clearest case is wall construction crossed with age.
\autoref{fig-cell-risk} shows the chance a facility has a first
confirmed release in a year, by tank age and wall type. For
single-walled facilities that annual chance climbs with age, from about
{[}X{]} percent in the youngest tanks to about {[}X{]} percent in the
oldest. For double-walled facilities it stays low and nearly flat, near
{[}X{]} percent, across the same ages. The two pull apart as tanks get
older. The lifetime picture is starker. Over a 30-year life about
{[}X{]} percent of single-walled facilities leak at least once, against
about {[}X{]} percent of double-walled facilities. Risk is not a
surprise that arrives at random. It is written in traits an underwriter
can read when the policy is written.

A single trait gradient is suggestive, but a premium rests on all the
traits together, and on whether the pattern holds for facilities the
model has not seen. So I estimate each facility's annual first-release
probability from its full set of observable traits, then check the
predictions out of sample. The model answers the underwriter's question.
For a facility that has never leaked, what is the chance of a first
release in the coming year? It uses only what an underwriter can see at
the start of the year. The traits are wall construction, tank age,
portfolio size, total capacity, and fuel type, with a state baseline. I
fit the model on pre-reform years only, and every facility's predicted
risk is scored by a model fit on data that left it out, so nothing here
is in-sample fit. Two facts show the predictions are usable. Risk
concentrates: the {[}10{]} percent of facilities the model ranks highest
hold {[}X{]} percent of next-year releases, about {[}Z{]} times what
random screening would catch (\autoref{fig-lift}). And the predictions
are right in levels, not just in rank: average predicted risk tracks the
observed release rate across the whole range of the score
(\autoref{fig-calibration}).

The model does more than rank facilities. It reproduces the level of
risk in each cell the market actually prices on. When I compare its
predicted release chance in each age-and-wall cell to the raw chance in
that cell, the points fall along the 45-degree line
(\autoref{fig-cell-gof}), so the model recovers the risk map rather than
just sorting facilities along it. It also separates leakers from
non-leakers well out of sample (ROC AUC {[}AUC{]}, precision-recall AUC
{[}PR\_AUC{]}). This predicted probability is what a risk-based premium
puts a price on. I reuse the same facility-level hazard as an input to
the structural model in §5. The lesson for the reform is plain. An
insurer can tell risky facilities from safe ones in advance, so a
premium can be made to rise with risk, and a flat per-tank fee chooses
not to.

\subsection{Pricing observable risk}\label{pricing-observable-risk}

After the fund closed, Texas facilities bought coverage from private
insurers, and I observe what they paid. I reconstruct each facility's
annual premium from the filed rate manuals of the state's largest UST
insurer, linked to the facility through the TCEQ contract records. The
schedule prices the same traits that predict releases. The per-tank
premium rises with tank age, and at any given age it is higher for
single-walled tanks than for double-walled tanks
(\autoref{fig-premium-by-age}). In the control states, where the public
fund still sets the price, I use the ASTSWMO survey of fund fees. Those
funds charge a flat amount per tank that does not move with age or wall
type.

The premium gradient is not arbitrary. It follows the risk map from
§4.1. When I line up each age-and-wall cell's premium against its
predicted release probability, the two move together: cells the model
flags as higher risk carry higher premiums
(\autoref{fig-actuarial-alignment}). The price an insurer charges rises
with the risk the same observables predict, which is what makes the
Texas schedule a risk-based price and not a relabeled flat fee.

Before the reform, Texas and the control states both paid a flat
per-tank fee, so their prices moved together. At the 1998 reform the
Texas price split off. \autoref{fig-premium-divergence} plots the
average per-tank premium in Texas against the control states from 1992
to 2020. The two track each other before 1998 and pull apart after, as
the Texas schedule starts to charge more for older single-walled tanks
while the control fee stays flat. Across matched facilities the premium
gap at the switch is about {[}\$X{]}. This change in the price of risk,
with the facilities themselves held fixed, is the variation the
difference-in-differences design uses in §4.3.

\subsubsection{The flat fee cross-subsidizes
risk}\label{the-flat-fee-cross-subsidizes-risk}

I can now compare what each facility pays under a flat fee to what its
risk actually costs. The fair premium for a facility is its leak risk
times the cost of a cleanup. I take the leak risk from the §4.1 hazard
model and the cleanup cost from incident-level claims in the states
where I have them, net of the deductible the facility would bear. The
claims show that cleanup cost barely depends on a tank's age or wall
type, so I use a single pooled cleanup cost across facilities and let
the whole risk gradient come from how often a facility leaks.\footnote{The
  leak frequency is the out-of-sample prediction from the §4.1 hazard
  model, a cross-validated penalized logistic regression on tank age,
  wall type, fuel, and capacity. The cleanup cost is the average claim
  paid net of the deductible, pooled across the states with usable
  claims data and deflated to 2023 dollars, about \$377,838 per release.}
The fair premium is what a facility's risk costs the fund each year, and
it is the vertical axis in the figures that follow. It runs from a few
hundred dollars for the safest facilities to over \$10,000 for the
riskiest in Tennessee, and the spread is similar in every state.

Today these facilities pay almost nothing toward that risk.
\autoref{fig-cross-paid} lines up each state's facilities from safest to
riskiest and marks what each one pays, its flat fund fee, against its
fair premium. The fee is a thin strip along the bottom. The share of its
fair premium a facility actually pays is small everywhere: 17 percent in
Tennessee, 14 percent in Louisiana, 1 percent in Colorado, and near zero
in New Mexico, which levies no flat fee and funds its program through a
per-gallon gas tax. The rest, from 83 percent in Tennessee to 100
percent in New Mexico, is covered by someone other than the facility
that carries the risk.

A flat fee that paid for itself would still cross-subsidize risk. Set
the flat fee at break-even, the average fair premium, so the fund
collects exactly what it pays out and needs no gas-tax help. In 2005
that fee is \$7,592 in Colorado, \$5,027 in New Mexico, \$4,160 in
Tennessee, and \$990 in Louisiana. \autoref{fig-cross-area} draws that
break-even line across the facilities. The safe facilities sit below it
and overpay, the blue area. The risky facilities sit above it and pay
less than they cost, the red area. The two areas are equal, so the safe
half's overpayments exactly fund the risky half's discount. The transfer
is 26 percent of premiums in New Mexico, 18 percent in Colorado, 18
percent in Louisiana, and 16 percent in Tennessee. The split of
facilities runs from 60 percent overpaying in Colorado to 77 percent in
New Mexico, where a small, very risky tail is carried by a large safe
majority. This is the cross-subsidy a flat fee creates, even when it is
priced to break even.

The cross-subsidy moves over time, and its two readings move in opposite
directions. \autoref{fig-cross-time} follows Tennessee across 2000,
2005, 2010, and 2015. Read in dollars, the transfer grows. The
break-even fee climbs as the fleet ages, because older tanks leak more
and the average cleanup cost rises, and a larger average pulls the gaps
between facilities up with it: the fee rises from \$3,845 to \$5,068 in
Tennessee, \$7,042 to \$8,728 in Colorado, and \$920 to \$1,137 in
Louisiana. Read as a share of premiums, the transfer shrinks, from 18 to
15 percent in Tennessee, 19 to 17 percent in Colorado, and 18 to 11
percent in Louisiana. What governs the share is not the level of risk
but its spread relative to the average, and that spread compresses as
the fleet ages together. New installations slow, the young low-risk end
thins out, facilities grow more alike, and there is proportionally less
to redistribute. The two readings answer different questions. The dollar
transfer is the size of the hidden subsidy a flat fee creates, and it is
rising, so the stakes of mispricing grow. The share is how distortionary
the fee is per dollar collected, and it is easing, because a maturing
fleet is a more homogeneous risk pool. Neither dominates, the flat fee
moves more money every year even as it becomes proportionally less
unfair. New Mexico is too thin to read a trend, so I report its 2005
level only.

A few choices shape these figures. Deductibles are large, so the fair
premium uses cleanup cost net of the deductible. New Mexico reports
gross cost, so I net out its deductible. The figures cover four states:
Colorado, Louisiana, New Mexico, and Tennessee. Louisiana has no
cleanup-cost data of its own, so it borrows the pooled average cost
across states. Pennsylvania is dropped because its hazard estimate is
degenerate. The cross-section is for 2005 unless noted, and the New
Mexico panel is small, about 70 facilities, so I use it for the 2005
level only and not for any trend.

\subsection{Design and identification}\label{design-and-identification}

I measure the effect of the reform on tank closure with a
difference-in-differences regression: \[
\text{closure}_{it}=\beta\,\text{did}_{it}+\gamma'M_{it}+\alpha_f+\delta_{ct}+\varepsilon_{it}.
\]

The outcome is one in the year tank \(i\) is permanently removed. The
term \(\text{did}_{it}\) turns on for Texas tanks from 1998, the year
the fund closed (December 22, 1998). The facility fixed effect
\(\alpha_f\) holds the site fixed, so the comparison is within a
facility over time. The cell-by-year fixed effect \(\delta_{ct}\) sets
the comparison group. A cell is a make and model: wall type, fuel,
capacity, and install cohort. So a Texas tank is compared only to
control tanks of the same make and model in the same year. The 1988
federal upgrade deadline is national, so it is common to treated and
control states and differences out of the difference-in-differences. The
cell-by-year fixed effects absorb cohort-specific mandate timing, and
the mandate indicators \(M_{it}\) are an additional and largely
redundant control, dropped in estimation for collinearity. No
installation cohort is dropped.\footnote{The federal deadline shares its
  December 1998 date with the Texas reform. That coincidence is
  precisely why the common-shock differencing of the
  difference-in-differences, not a cohort drop, is what isolates the
  Texas risk-based-pricing effect.} The coefficient \(\beta\) is the
average change in a Texas tank's annual closure probability after the
reform, within its make-model-cohort and year. This \(\beta\) is causal
under one assumption. Without the reform, Texas tanks and their matched
controls in the same cell would have followed the same closure path. The
event study in the next subsection tests it. The headline sample is 9.76
million tank-years across 117,250 facilities in 18 states.

To see the timing, I replace the single post indicator with one
coefficient per year relative to the reform: \[
\text{closure}_{it}=\sum_{k\ne -1}\beta_k\,\mathbf 1\{k_{it}=k\}\,\text{TX}_f+\gamma'M_{it}+\alpha_f+\delta_{ct}+\varepsilon_{it},
\]

where \(k=t-1998\) is event time and I normalize to the year before the
reform, \(k=-1\). The pre-reform coefficients are the parallel-trends
test: if Texas and control tanks moved together before the reform, the
\(\beta_k\) for \(k<0\) sit near zero. The post-reform coefficients
trace the path of the response.

A firm decides at the facility, not tank by tank, so I carry the same
design up to the portfolio. I build each facility's no-reform benchmark
from how untreated tanks of the same make and model behaved,
\(\hat Y^0_{ft}=\sum_c s_{fct}\,\hat\delta_{ct}\), where \(s_{fct}\) is
the facility's share of tanks in cell \(c\) and \(\hat\delta_{ct}\) is
the untreated closure rate in that cell. The facility regression mirrors
the tank one: \[
Y_{ft}=\beta\,\text{did}_{ft}+\mu\,\hat Y^0_{ft}+\gamma'M_{ft}+\alpha_f+\delta^{\text{mix}}_{m(f),t}+\varepsilon_{ft}.
\]

It holds the facility fixed with \(\alpha_f\), holds a
portfolio-mix-by-year fixed effect \(\delta^{\text{mix}}_{m(f),t}\)
fixed (the facility's tank mix interacted with the year, the facility
twin of the tank make-model-by-year cell), and keeps \(\hat Y^0_{ft}\)
as a control. The outcome \(Y_{ft}\) is any of the portfolio margins. As
a check I also regress the imputation residual
\(\tau_{ft}=\overline{\text{closure}}_{ft}-\hat Y^0_{ft}\) on the
treatment, which gives the same answer.

I cluster standard errors by state. Texas is the only treated state, so
I compute valid standard errors with a wild cluster bootstrap.

\subsection{Does the policy work?}\label{does-the-policy-work}

Risk-based pricing raised tank closures. After the reform, the annual
chance that a Texas tank is permanently removed rose by 1.58 percentage
points (s.e. 0.40), on a pre-reform control mean of about 2 percent
(\autoref{tbl-stepped-did}). That is close to an 80 percent increase,
measured within make-model-cohort cells and holding the facility fixed.
The reform moved real capital.

The response lines up with the reform date. \autoref{fig-es-tank} plots
the year-by-year closure gap between Texas and control tanks. Before the
reform the gap is flat and close to zero, so the two groups were on
parallel paths. At the reform the gap opens, and it stays open for years
after. The effect is a clean break at 1998, not the continuation of a
pre-existing trend.

The effect is the reform, not something else that hit Texas in 1998. A
recession or an oil-price swing would close tanks across the board. The
reform instead repriced risk, so it should bite hardest on the tanks
whose price rose most, the old single-walled ones, and barely touch new
double-walled tanks. To test this I set aside the make-model fixed
effect, which would otherwise absorb wall and age, and interact the
treatment with tank age and wall type. The closure response is {[}β{]}
points larger for old single-walled tanks and near zero for new
double-walled tanks (\autoref{tbl-risk-conc}). It lands on the tanks
whose premium rose most under risk-based pricing, which a uniform
statewide shock cannot produce.

The same result appears when I move from the tank to the facility. With
the facility fixed and the portfolio-mix-by-year baseline held fixed,
the share of a facility's tanks that close rose by 1.62 percentage
points after the reform (s.e. 0.46), and the chance a facility closed
any tank rose by 1.83 points (s.e. 0.46) (\autoref{tbl-fac-att}). A
second route, which imputes each facility's no-reform closure share and
regresses the residual on the treatment, gives a similar 1.18 points
(s.e. 0.38).\footnote{Route A is the direct fixed-effects estimate on
  the portfolio-mix-by-year cell. Route B imputes the no-reform baseline
  from untreated tanks of the same make and model and uses a plain
  facility-and-year fixed effect. The mix-by-year cell is what lifts
  Route A above Route B. The two bracket the facility-level effect.} So
the reform changed behavior at the level a firm actually decides, the
whole facility, not only the single tank.

The facility path tells the same story. \autoref{fig-es-fac} plots the
facility-level closure gap by year, with the mix-by-year baseline and
the composition control held fixed. It is flat before the reform and
rises after, matching the tank-level path. The portfolio response and
the tank response line up.

The reform also lowered leaks. Annual confirmed leak discovery in Texas
fell by about 1.2 percentage points (s.e. 0.003) after the reform, on a
pre-reform control mean of about 1.5 percent (\autoref{tbl-leak-did}).
This drop is the environmental payoff, and it is not a mechanical
artifact of closing more tanks. Closing a tank usually triggers an
inspection that can surface a hidden leak, which would push discovery
up, yet discovery fell. It fell through both the routine-monitoring
channel and the closure-inspection channel (\autoref{tbl-leak-decomp}),
so the reform lowered the true rate of leaks, not just how they are
found. The leak event study is noisier than the closure one, so I read
this result from the average effect rather than the year-by-year path.

\subsection{Who, where, and what}\label{who-where-and-what}

Facing the new prices, a facility had several moves on its portfolio,
and I line them up with the actions in the structural model. It can
downsize, cutting both tank count and capacity. It can consolidate,
cutting tanks while holding capacity. It can reconfigure up, cutting
tanks while adding capacity, which is rare. It can replace tanks at the
same capacity, or it can exit. I split the contraction moves by net
gallons, so shedding capacity is kept apart from holding it, and I
estimate each margin on all facilities, not only the ones that close.
The reform raised genuine downsizing by 0.26 points (s.e. 0.09) and the
rare capacity-adding reconfiguration by 0.16 points (s.e. 0.03), while
same-capacity consolidation did not move, at −0.04 points and not
significant. The contraction the reform produced is real capacity
shedding, not a relabeling of the same fleet.

I also trace each contraction margin year by year. The downsizing path
wobbles before the reform, so I read it as descriptive
(\autoref{fig-es-downsize}). The consolidation path shows no clear break
(\autoref{fig-es-consolidate}), matching its null average. The
replacement path is moderate with an acceptable pre-period
(\autoref{fig-es-replace}).

The response is concentrated in motor-fuel retailers. At the tank level,
facilities whose tanks are all motor fuel close about 4.3 percentage
points more than other facilities after the reform (p \textless{}
0.001). Their total response is near 5 percentage points. Having any
gasoline tank, which describes most facilities, does not separate the
response. That differential is about −0.5 points and is not significant.
At the facility level the pattern holds and runs across the whole
portfolio. Gas stations close about 1.8 percentage points more than
other facilities (s.e. 0.16), and the gas-station differential is
positive on every margin I measure, including exit and replacement
(\autoref{tbl-hte-tank}, \autoref{tbl-fac-hte}). Motor-fuel retailers
carry the response.

The response also sorts by how big a facility is and how old its tanks
are, and it sorts the same way on both. I interact the treatment with
the facility's total-capacity bin, four bins from small to large, and
with its tank vintage at the reform. Small facilities exit. The smallest
bin, under 9,000 gallons, raises its closure share by 6.3 points (s.e.
0.5) and its exit rate by 6.6 points (s.e. 0.5), and that response falls
toward zero across the larger bins (\autoref{tbl-size-hte}). Large
facilities reinvest instead. The largest bin, over 30,000 gallons,
barely exits but downsizes 1.0 point and replaces tanks 0.7 points, each
rising monotonically across the bins (p \textless{} 0.001 for every
differential). Vintage sorts the same way. The oldest tanks, installed
before 1975, exit 2.7 points more than the youngest cohort (s.e. 0.3)
and do not reinvest, while the youngest cohort, 1989 to 1998, barely
exits but downsizes 0.5 points (s.e. 0.04) and replaces tanks 0.7 points
(s.e. 0.04) (\autoref{tbl-vintage-hte}). The marginal facilities, small
and old, leave the market. The viable ones, large and young, upgrade and
stay.

Where a facility sits matters too. I interact the treatment with
location and market measures fixed at the reform: rural, low-population,
low-income, and high-poverty areas from the 2000 Census, and a
thin-market indicator for gas retailers with few active competitors
within a mile in 1998. The competition measure is set before the reform,
so it does not pick up the reform's own effect on the market. Two
patterns stand out. Where demand is captive, facilities stay open. Rural
facilities exit 0.42 points less than others (s.e. 0.03), and
thin-market gas retailers exit 1.64 points less (s.e. 0.09). Where
facilities serve poor areas, they exit but do not reinvest. Low-income
and high-poverty facilities exit more, by 0.22 and 0.18 points (s.e.
0.04 each), yet downsize and replace less, with low-income facilities
downsizing 0.86 points less (s.e. 0.02) and replacing 0.32 points less
(s.e. 0.01). So the policy retired tanks fastest in poor areas, with no
upgrading there, and slowest where a single station holds a captive
market.

Put together, the reduced form answers the policy question. Risk-based
pricing made facilities retire risky tanks, the response was a clean
break at the reform, and it lowered leaks. It fell hardest on the
operators carrying the most risk, with facilities upgrading as well as
exiting. Because the response concentrated in the oldest single-walled
tanks, where the premium rose most, rather than spreading evenly (§4.4),
a statewide shock that happened to coincide with 1998 cannot explain it.
What the reduced form cannot do is put these responses in dollars, price
the policies Texas never ran, or separate deterrence from selection.
That is the work of the structural model in §5.

\section{Dynamic Structural Model}\label{dynamic-structural-model}

The reduced form shows that the 1998 reform changed closure behavior. It
cannot value that change in dollars, price a policy Texas never ran, or
say how far the reform moved firms toward the social optimum. To answer
those questions I build a dynamic structural model of how a facility
manages its tank portfolio. Each year a facility chooses what to do with
its tanks. It can keep operating them, remove some, replace some, or
shut down. It weighs the operating value of its tanks against the cost
of the risk they carry and the cost of acting. I estimate the model on
Texas and the uniform-pricing control states, then re-solve it under
policies that were never tried.

The agent is the facility, not the tank. A facility with several aging
tanks does not decide tank by tank. It decides how many of its tanks to
remove as one choice. The state has to track the facility's whole
portfolio, not one representative tank. A model built on a single tank
cannot tell a facility that slowly retires an old fleet apart from one
that replaces tanks one for one, and those two paths carry different
consequences for a removal mandate. Tracking the portfolio is what lets
the model price the partial-closure margin the reduced form found.

\subsection{Measured operating revenue (first
stage)}\label{measured-operating-revenue-first-stage}

The model needs to know what a facility earns, not only what it pays for
insurance. I build a per-period operating-revenue measure for each
facility from public data. The retail gasoline price and gallons sold by
state and year come from the EIA State Energy Data System. The state
gasoline tax comes from the FHWA ``Tax Rates by Motor Fuel and State''
series. The wholesale price refiners charge resellers comes from the EIA
refiner ``sales for resale'' prices, reported by region (PADD). The
federal gasoline tax is 18.4 cents a gallon. Facility size comes from
the tank registry: I sum each facility's tank capacity and sort
facilities into four size groups. The series cover the 18 study states
from 1999 to 2020.

Most of the retail price of gasoline is the cost of the fuel itself, so
the retail price is a poor measure of what a station earns. The right
measure is the margin: the retail price minus the wholesale price, the
state tax, and the federal tax. The margin is the money a station makes
on a gallon. A facility's revenue is its margin times how much fuel it
moves, and throughput scales with size. So I measure revenue as capacity
times the local fuel margin. Everything is in nominal dollars, to match
the nominal premium and deductible already in the model. I average the
margin within three eras rather than tracking every year, so a facility
forecasts a steady local margin instead of chasing the year-to-year
swing.

The model used to carry one free number per size group for the operating
payoff. Now that size profile is tied to measured revenue. The term
becomes \(\psi R\), where \(R\) is the measured revenue and \(\psi\) is
a single weight I estimate. \(\psi\) says how much a dollar of revenue
moves a facility's choices, on the same footing as the weights on the
premium and the deductible. This follows Sweeting (2013): measure
revenue outside the dynamic model, then estimate only the weight. A
constant linking storage capacity to gallons sold is folded into
\(\psi\), so \(\psi\) is a weight, not a dollar amount. Because revenue
enters the post-action payoff, a facility that downsizes or exits gives
up measured revenue right away, which ties the size margins to real
dollars. \(\psi\) is identified mainly across facilities, from a
facility's size crossed with its state's fuel margin, not from the
national time path of prices.\footnote{One state's early-period margin
  is bounded, not measured. North Carolina sits in the Lower Atlantic,
  which the Gulf supplies through the Colonial Pipeline, but EIA reports
  wholesale prices only for the whole East Coast region, which is pulled
  up by the Northeast. With no sub-region wholesale price published,
  North Carolina's margin before 2014 is floored at a small positive
  value. It is one cell of fifty-four, and the state fixed effect
  absorbs the rest.}

\subsection{State space}\label{state-space}

The portfolio \(c\) records how many tanks a facility holds across 16
categories. The categories cross two wall types, single-walled and
double-walled, with eight five-year age bands. A facility holds between
one and six tanks. Counting every combination of tank counts across the
categories gives 74,612 distinct portfolios. A separate term \(G\) sorts
each facility into one of four bins by total storage capacity. I keep
capacity apart from the portfolio because capacity drives revenue, while
tank count and age drive the premium and the hazard. Two facilities with
the same tanks can still differ in scale. Crossing portfolios with
capacity bins gives a state space of \(74{,}612 \times 4 = 298{,}448\)
states.

Two facilities show what a state is. The first holds three tanks: two
single-walled tanks in the oldest band and one double-walled tank in the
youngest, with capacity in the smallest bin. The second holds four
double-walled tanks in a middle band, with capacity in the largest bin.
The two face different premiums and different hazard even under the same
contract, because premium and hazard depend on the full portfolio, not
on a summary like tank count or average age. All facility heterogeneity
in the model is where the facility sits in this state space.

Most facilities sit in a simple part of this space. About 82 percent of
facility-years hold tanks of a single wall type and age band, so the
portfolio collapses to a count for them. The full richness of the state
is used only by the large, mixed facilities, which are also the ones
whose partial-closure choices the reduced form flagged.

\subsection{Actions}\label{actions}

Each year a facility chooses an action \(a = (k, m)\). It removes \(k\)
tanks and installs \(m\) new ones. Removal follows a fixed order: oldest
tanks first, single-walled before double-walled, smallest first. That
order matches the order in 84 percent of removal events in the data, so
I build it in as a rule rather than estimate a separate choice over
which tank leaves. Installation is restricted to double-walled tanks,
because federal and state rules ended new single-walled installation
over the sample.\footnote{New single-walled installation was foreclosed
  by state requirements that new tanks and piping be secondarily
  contained, that is, double-walled. Most states followed the federal
  secondary-containment requirement added by the Energy Policy Act of
  2005. Effective dates run from 1989 in Massachusetts and 1991 in Maine
  through 2013 in Kansas, with most concentrated in 2007 to 2009.} There
is no separate choice over wall type when a facility installs.

The menu covers the behaviors the reduced form measured. Maintaining is
\((0,0)\). Downsizing is \((k,0)\), removing tanks and operating with
fewer. Replacing or upgrading is \((k, m>0)\), removing tanks and
installing new double-walled ones. A facility can also exit, which is
absorbing. Removal and installation are one joint choice, not two steps,
so a facility's plan to replace tanks can shape how many it removes in
the same year. A two-step model would force a facility to decide whether
to close before it could decide whether to reopen, which is not how
replacement looks in the data.

Premium and hazard are evaluated on the fleet left after the action,
before that year's aging. A facility that downsizes sees its premium
fall in the same year it removes tanks, not later. This matches how an
insurer re-rates a facility after its tank count changes.

\subsection{Flow utility}\label{flow-utility}

A facility that does not exit gets a per-period payoff from operating
the fleet \(n'\) left after its action: \[
u_a(s,e)=\phi_G+\alpha_e-\big(\gamma_{\text{pool}}+\gamma_{\text{RB}}\,\mathbf 1[\text{RB}_e]\big)\,E(n',e)-k\,c_{\text{rem}}-m\,c_{\text{inst}}.
\] Here \(\phi_G\) is the operating value of the fleet, allowed to vary
with the capacity bin \(G\). \(\alpha_e\) is a fixed effect for the
facility's regulatory environment. \(E(n',e)\) is the total risk
exposure the fleet carries. \(c_{\text{rem}}\) and \(c_{\text{inst}}\)
are the per-tank costs of removing and installing.
\(\gamma_{\text{pool}}\) and \(\gamma_{\text{RB}}\) convert exposure
into the same units as operating value, and I return to them below.

Total risk exposure is the sum of two costs a facility bears for holding
risky tanks: \[
E(n',e)=P(n',e)+H(n')\,D_e.
\] \(P(n',e)\) is the insurance premium the facility pays on the fleet
\(n'\) under its contract \(e\). \(H(n')\) is the facility's hazard, the
chance that any of its tanks leaks in the year. \(D_e\) is the
deductible the facility pays out of pocket when a leak occurs, so the
second term is the expected out-of-pocket cost of a leak. Both pieces
rise with older and single-walled tanks. The hazard \(H\) comes from the
facility-level model in §4.1, the premium \(P\) from the reconstructed
schedule in §4.2, and the deductible \(D_e\) from the contract terms.
All three are first-stage inputs, estimated outside the dynamic model
and held fixed.

The model puts one coefficient \(\gamma_{\text{pool}}\) on total
exposure, not separate coefficients on the premium and on the expected
out-of-pocket cost. A dollar of premium and a dollar of expected
out-of-pocket loss enter the payoff the same way. I do this because the
data cannot tell the two apart. Under risk-based pricing the premium
moves with the hazard, so the premium and the expected loss are
collinear. The premium also varies mostly across states, which the
environment fixed effects absorb. And the Texas premium data begin only
in 2006, so there is no clean within-Texas window in which to separate
them. I report the pooled response to total exposure and treat the
premium-versus-out-of-pocket split as not identified in this data. That
non-identification is itself a result. It says a facility responds to
the size of its risk bill, not to which channel delivers it.

The second coefficient \(\gamma_{\text{RB}}\) is an extra response that
applies only under risk-based pricing. It lets a dollar of exposure move
behavior more, or less, when the price is risk-based than when it is a
flat fund fee. A positive \(\gamma_{\text{RB}}\) means firms respond
more strongly to exposure under risk-based pricing, which is the sign
that the pricing mechanism does work beyond the level of the bill. A
value near zero means a dollar is a dollar whatever the regime. In this
sample risk-based pricing means Texas, so \(\gamma_{\text{RB}}\) also
picks up anything specific to Texas. I read it as suggestive of the
mechanism, not as a clean causal estimate, and I say so again where it
matters for the counterfactuals.

The environment fixed effects \(\alpha_e\) are load-bearing. They absorb
everything that differs across states for reasons unrelated to exposure:
geology, enforcement, the local economy, data coverage. Holding them
fixed is what leaves a clean within-state exposure signal for
\(\gamma_{\text{pool}}\) to read. The operating value \(\phi_G\) is the
payoff from running a fleet at a given scale. In the estimated model I
replace the free \(\phi_G\) intercepts with the measured-revenue term
from §5.1, so operating value is \(\psi\,R\), where \(R\) is the
measured revenue of a facility of that size in that state and era and
\(\psi\) is the single weight that scales it. This ties operating value
to measured revenue instead of free intercepts, which came back nearly
flat across capacity bins.

A facility that exits gets a one-time terminal value: \[
u_{\text{exit}}(s)=\kappa_1\,N(s).
\] \(N(s)\) is the number of tanks the facility holds. \(\kappa_1\)
measures how exit value scales with size. I normalize the constant part
of exit value to zero, which pins the level of the value function, so
\(\kappa_1\) measures only how exit value grows with the number of
tanks.

\subsection{Bellman equation and
estimation}\label{bellman-equation-and-estimation}

Facilities are forward-looking. The value of a state solves \[
V(s)=\max\Big\{\max_{(k,m)}\big[u_a(s,e)+\beta\,\mathbb E[V(s')\mid s,a]\big],\;u_{\text{exit}}(s)\Big\}.
\] The discount factor \(\beta=0.9957\) is fixed in advance rather than
estimated. Exit is absorbing. Each action carries an independent type-1
extreme value shock, which gives closed-form logit choice probabilities
given the state.

Estimation runs in two stages. The first stage builds the premium
schedule, the hazard model, the deductibles, the capacity transitions,
and the revenue measure directly from the claims, premium, and registry
data, and holds them fixed. This is the standard two-stage approach in
this literature. It lets the data on prices and risk be taken as given
before I ask what cost and preference parameters fit the closure
behavior observed under those prices.

The second stage estimates
\(\theta=(\psi,\gamma_{\text{pool}},\gamma_{\text{RB}},c_{\text{rem}},c_{\text{inst}},\kappa_1)\),
with the environment fixed effects \(\alpha_e\) estimated alongside. One
parameter vector \(\theta\) is shared across all environments: 14
uniform-pricing control states and three Texas risk-based rate eras, 17
in all, with Texas the reference. Sharing one \(\theta\) across
environments that price risk differently is the main source of
identifying variation, which the next section takes up.

I estimate \(\theta\) by nested pseudo-likelihood
\citep{aguirregabiria_sequential_2007}. The method avoids solving the
full dynamic program at every trial value of \(\theta\). It starts from
the choice probabilities seen in the data, holds them fixed to turn the
dynamic problem into a static one in which the value function is linear
in \(\theta\), estimates \(\theta\) by maximizing the pseudo-likelihood,
then recomputes the model's implied choice probabilities and repeats
until \(\theta\) and the probabilities settle together. At 298,448
states, solving the value function from scratch inside an optimizer
would not be feasible. With the choice probabilities fixed the value
function solves from the linear inversion \(V=(I-\beta M)^{-1}R\), which
I compute with an iterative solver (BiCGSTAB) preconditioned on the
tank-aging structure rather than by direct inversion. The estimating
sample is aggregated facility-year choice counts, Texas 2006 to 2020 and
the control states 1999 to 2020. The 2006 start for Texas is forced by
when the premium data begin.

\subsection{Identification}\label{identification}

Identification of \(\theta\) rests on three things. The first is a set
of modeling assumptions grounded in patterns in the data. The second is
the set of assumptions that comes with choosing a dynamic discrete
choice model estimated by NPL. The third is the actual comparison in the
data that fixes the size of each parameter once the first two hold. I
take them in turn, and I am explicit that some parameters rest on
stronger variation than others.

\subsubsection{Assumptions grounded in the
data}\label{assumptions-grounded-in-the-data}

The removal order, oldest first and single-walled before double-walled,
is the order in 84 percent of removal events, so I build it in as a
rule. Installation is double-walled only because the single-walled
install ban removed the alternative. The premium schedule, hazard,
deductibles, capacity transitions, and revenue are built from observed
records and held fixed, because they are directly observed. The cost and
preference parameters in \(\theta\) are not observed and must be
estimated.

\subsubsection{Assumptions from the model
class}\label{assumptions-from-the-model-class}

These are standard for the model class, not special to tanks. Each
action's shock is an independent draw from a type-1 extreme value
distribution, which gives the logit choice probabilities. Facilities
have rational expectations: each knows its contract and correctly
anticipates how its own portfolio evolves. The NPL mapping satisfies the
regularity conditions needed for the iteration to converge to the
parameters that generated the data. Convergence diagnostics are reported
with the fit.

\subsubsection{What pins each parameter}\label{what-pins-each-parameter}

Conditional on those assumptions, each parameter is pinned by a distinct
comparison. \(\gamma_{\text{pool}}\), the response to total exposure, is
identified from within-state variation in exposure. At a fixed
environment, facilities with older or single-walled tanks, or facing
higher deductibles, carry more exposure and downsize or exit more. The
environment fixed effects are what make this clean, because they strip
out cross-state confounds and leave within-state variation. The age-free
levers that discipline it are wall type and the deductible.
Single-walled versus double-walled at the same age changes hazard, and
the deductible changes out-of-pocket cost without being driven by age,
so both move exposure for reasons other than a tank getting older. The
premium-versus-out-of-pocket split is not identified: with separate
coefficients the premium coefficient collapses to near zero, because
premium moves with hazard under risk-based pricing, because premium
varies mostly across states that the fixed effects absorb, and because
Texas premiums begin only in 2006. I therefore report the pooled
response. \(\gamma_{\text{RB}}\), the extra response under risk-based
pricing, is identified from the difference in the exposure-response
slope between the uniform control states and risk-based Texas. Because
risk-based pricing is Texas in this sample, this comparison conflates
the mechanism with Texas-specific factors, so it is suggestive, not a
clean estimate of the mechanism. \(\psi\), the weight on operating
revenue, is identified across facilities, from a facility's size crossed
with its state's fuel margin (§5.1), not from the national time path of
prices. \(c_{\text{rem}}\) and \(c_{\text{inst}}\) are identified from
how often facilities downsize and replace, and from the split among
closing facilities between downsizing, replacement, and exit.
\(\kappa_1\) is identified from how the full-exit rate varies with the
number of tanks, net of operating value.

This strategy has clear limits. The premium-versus-out-of-pocket split
is not identified, so I report the pooled response.
\(\gamma_{\text{RB}}\) conflates risk-based pricing with Texas. The
hazard model is only as precise as the leak data. The removal rule
matches 84 percent of events, leaving a small share it does not.

\section{Estimates, Counterfactuals, and
Welfare}\label{estimates-counterfactuals-and-welfare}

\subsection{Estimates}\label{estimates}

\autoref{tbl-struct-estimates} reports the estimated parameters. The
headline pooled coefficients, \(\gamma_{\text{pool}}\) on total exposure
and \(\gamma_{\text{RB}}\) on the risk-based regime, are {[}pending the
converged pooled run{]}. The anchor in hand is the earlier
two-coefficient fit, in which the single exposure coefficient is 4.32,
the removal cost is 3.85, the install cost is 2.04, and the exit-size
gradient \(\kappa_1\) is 0.93, with a log-likelihood of \(-291{,}648.5\)
across 17 environments. Operating value came back nearly flat across
capacity bins under free intercepts, which is part of why I tie it to
measured revenue. All values are in model units of \$10,000 per
facility-year. Standard errors are {[}pending{]}. The significance of
\(\gamma_{\text{RB}}\) is what the regime reading turns on.

{[}Pending the pooled fit: the reading of the headline. A positive
\(\gamma_{\text{pool}}\) means more exposure pushes facilities off
maintaining toward downsizing and exit. The sign and size of
\(\gamma_{\text{RB}}\) is the regime-mechanism result, whether a dollar
of exposure bites harder under risk-based pricing.{]}

\subsection{Counterfactuals}\label{counterfactuals}

With \(\theta\) estimated, I re-solve the model under policies Texas
never ran. Each counterfactual changes one term in the payoff and leaves
the rest of the model unchanged, then re-solves the value function and
the resulting fleet. I keep the environment fixed effects in the
counterfactual, because they stand for real and persistent features of
each state rather than a temporary effect, so the baseline reproduces
observed behavior and each policy is a clean change against it. Part of
each fixed effect may reflect data coverage rather than behavior, which
I acknowledge in reading the results. Welfare is computed on the
observed distribution of facilities, with no reweighting.

The first counterfactual asks what Texas would have done if it had
stayed on uniform pricing. I replace the Texas risk-based contract, both
its premium and its deductible, with a representative uniform-state
contract and re-solve. I report it as a decomposition. The first piece
changes only the contract and keeps the risk-based response. The second
piece changes the contract and also turns off the regime-extra response
\(\gamma_{\text{RB}}\). The difference between the two isolates the
risk-based-pricing mechanism, separate from the change in the dollar
bill. This is the headline risk-based-versus-uniform experiment.

The second counterfactual subsidizes new double-walled tanks by lowering
the install cost, and sweeps the subsidy across levels. I report how the
replacement rate and welfare respond, and how the response differs
between the risk-based and uniform regimes. The subsidy adds a
government outlay to the welfare account.

The third counterfactual forces removal of single-walled tanks past age
20. A facility can comply by downsizing or by upgrading to double-walled
tanks. The mandate trades the firm's cost of forced early action against
the leak damage it averts. I compare its welfare to the gradual response
the price instruments produce.

\subsection{Welfare accounting}\label{welfare-accounting}

I score each counterfactual on one identity. Net social welfare is
producer surplus minus external leak damage minus government outlay, all
in present value. Producer surplus is expected discounted facility
value. External damage is expected discounted leaks times the external
cost per release \(E\). Government outlay is the subsidy, which enters
only in the subsidy counterfactual. Premiums are a transfer inside firm
value, so they net out of social welfare. The external cost \(E\) is a
researcher input, set at a placeholder of \$50,000 per release. This
identity makes the wedge explicit. A facility pays only its deductible
out of pocket, the insurance pool covers the rest of the cleanup, and
the health externality \(E\) falls outside both. So even a fairly priced
risk-based premium internalizes only the insured part of the social
cost, and the gap to the first best is \(E\).

\section{Conclusion}\label{conclusion}

\section*{Tables and Figures}\label{tables-and-figures}
\addcontentsline{toc}{section}{Tables and Figures}

\subsection{---- Section 3: data and descriptives
----}\label{section-3-data-and-descriptives--}

\clearpage
\null\vfill
\begin{center}
\phantomsection\captionof{table}{Baseline characteristics of study facilities at the reform date (December 22, 1998)}\label{tbl-summary}
\vspace{6pt}
\fbox{\textit{[pending table: Output/Tables/T\_Baseline\_Characteristics\_Slide]}}
\par\vspace{10pt}
\begin{minipage}{0.85\textwidth}\small\textit{Notes:} Texas and the 17 control states in separate columns. Alive-at-reform means at least one tank active on December 22, 1998. Capacity is the median, which is robust to large-facility outliers.\end{minipage}
\end{center}
\vfill\clearpage

\clearpage
\null\vfill
\begin{center}
\phantomsection\captionof{figure}{Tank count per incumbent facility at the reform date}\label{fig-baseline-tanks}
\par\vspace{8pt}
\fbox{\textit{[pending figure: Output/Figures/Figure\_Baseline\_TankCount]}}
\par\vspace{8pt}
\begin{minipage}{0.85\textwidth}\small\textit{Notes:} Active tank count per facility on December 22, 1998. Control states hold more single-tank facilities; Texas is shifted toward three- and four-tank sites.\end{minipage}
\end{center}
\vfill\clearpage

\clearpage
\null\vfill
\begin{center}
\phantomsection\captionof{figure}{Total tank capacity per incumbent facility at the reform date}\label{fig-baseline-capacity}
\par\vspace{8pt}
\fbox{\textit{[pending figure: Output/Figures/Figure\_Baseline\_Capacity]}}
\par\vspace{8pt}
\begin{minipage}{0.85\textwidth}\small\textit{Notes:} Total storage capacity per facility on December 22, 1998, clipped at 60,000 gallons. Control states carry a large mass of small single-tank facilities; Texas capacity is shifted right.\end{minipage}
\end{center}
\vfill\clearpage

\clearpage
\null\vfill
\begin{center}
\phantomsection\captionof{figure}{Distribution of cleanup costs per confirmed release}\label{fig-cost-distribution}
\par\vspace{8pt}
\fbox{\textit{[pending figure: Output/Figures/Figure\_Cost\_Distribution]}}
\par\vspace{8pt}
\begin{minipage}{0.85\textwidth}\small\textit{Notes:} Cleanup cost per LUST site across the six control states with usable trust-fund claims data. The distribution has a long right tail.\end{minipage}
\end{center}
\vfill\clearpage

\subsection{---- Section 4.1: observable risk
----}\label{section-4.1-observable-risk--}

\clearpage
\null\vfill
\begin{center}
\phantomsection\captionof{figure}{Annual first-release probability by tank age and wall type}\label{fig-cell-risk}
\par\vspace{8pt}
\fbox{\textit{[pending figure: Output/Figures/Figure\_CV\_CellRisk]}}
\par\vspace{8pt}
\begin{minipage}{0.85\textwidth}\small\textit{Notes:} Annual probability of a first confirmed release by tank age band and wall type. The single-walled gradient rises with age; the double-walled line stays low and flat.\end{minipage}
\end{center}
\vfill\clearpage

\clearpage
\null\vfill
\begin{center}
\phantomsection\captionof{figure}{Risk concentration in the highest-ranked facilities}\label{fig-lift}
\par\vspace{8pt}
\fbox{\textit{[pending figure: Output/Figures/FigureA\_CV\_Lift\_State]}}
\par\vspace{8pt}
\begin{minipage}{0.85\textwidth}\small\textit{Notes:} Out-of-sample lift. Facilities are ranked by predicted first-release probability; the curve plots the share of next-year releases captured against the share of facilities screened.\end{minipage}
\end{center}
\vfill\clearpage

\clearpage
\null\vfill
\begin{center}
\phantomsection\captionof{figure}{Calibration of predicted first-release probability}\label{fig-calibration}
\par\vspace{8pt}
\fbox{\textit{[pending figure: Output/Figures/Figure\_CV\_Calibration]}}
\par\vspace{8pt}
\begin{minipage}{0.85\textwidth}\small\textit{Notes:} Out-of-sample calibration. Mean predicted risk against the observed release rate across bins of the predicted score. Points on the 45-degree line indicate predictions correct in levels.\end{minipage}
\end{center}
\vfill\clearpage

\clearpage
\null\vfill
\begin{center}
\phantomsection\captionof{figure}{Predicted versus observed release probability, by age-and-wall cell}\label{fig-cell-gof}
\par\vspace{8pt}
\fbox{\textit{[pending figure: Output/Figures/Figure\_CV\_CellRisk\_GoF]}}
\par\vspace{8pt}
\begin{minipage}{0.85\textwidth}\small\textit{Notes:} Out-of-sample fit at the cell level. Predicted release probability against the raw rate in each tank age-and-wall cell. Points on the 45-degree line indicate the model recovers the cell-level risk map.\end{minipage}
\end{center}
\vfill\clearpage

\subsection{---- Section 4.2: pricing observable risk
----}\label{section-4.2-pricing-observable-risk--}

\clearpage
\null\vfill
\begin{center}
\phantomsection\captionof{figure}{Total facility premium by tank age band, Texas versus control states}\label{fig-premium-by-age}
\par\vspace{8pt}
\fbox{\textit{[pending figure: Output/Figures/AE\_X3\_Premium\_by\_Age]}}
\par\vspace{8pt}
\begin{minipage}{0.85\textwidth}\small\textit{Notes:} Texas points are the mean reconstructed premium by tank age band; control points are the state flat fee times tank count. The Texas premium rises with age and is higher for single-walled tanks; the control fee is flat.\end{minipage}
\end{center}
\vfill\clearpage

\clearpage
\null\vfill
\begin{center}
\phantomsection\captionof{figure}{Texas premiums against predicted release probability, by risk cell}\label{fig-actuarial-alignment}
\par\vspace{8pt}
\fbox{\textit{[pending figure: Output/Figures/Figure\_Actuarial\_Alignment]}}
\par\vspace{8pt}
\begin{minipage}{0.85\textwidth}\small\textit{Notes:} Mean Texas premium against the predicted first-release probability for each age-and-wall cell. The two move together, so the premium gradient tracks predicted risk.\end{minipage}
\end{center}
\vfill\clearpage

\clearpage
\null\vfill
\begin{center}
\phantomsection\captionof{figure}{Average per-tank premium by year, Texas versus control states}\label{fig-premium-divergence}
\par\vspace{8pt}
\fbox{\textit{[pending figure: Output/Figures/04e\_Premium\_by\_Year\_Treatment]}}
\par\vspace{8pt}
\begin{minipage}{0.85\textwidth}\small\textit{Notes:} Average per-tank premium in Texas and the control states, 1992 to 2020. The series track before 1998 and pull apart after, as the Texas schedule begins to charge more for older single-walled tanks.\end{minipage}
\end{center}
\vfill\clearpage

\subsection{---- Section 4.2.1: the cross-subsidy
----}\label{section-4.2.1-the-cross-subsidy--}

\clearpage
\null\vfill
\begin{center}
\phantomsection\captionof{figure}{What facilities pay versus their fair premium, by state}\label{fig-cross-paid}
\par\vspace{8pt}
\fbox{\textit{[pending figure: Output/Figures/Fig\_CrossSub\_Paid\_Panel]}}
\par\vspace{8pt}
\begin{minipage}{0.85\textwidth}\small\textit{Notes:} Facilities in each state ordered from safest to riskiest. The fair premium (leak risk times cleanup cost, net of the deductible) is the full height; the flat fund fee actually paid is the lower band. The share of the fair premium a facility pays is 17 percent in Tennessee, 14 percent in Louisiana, 1 percent in Colorado, and near zero in New Mexico, which charges no flat fee and funds its program through a per-gallon gas tax. 2005 cross-section.\end{minipage}
\end{center}
\vfill\clearpage

\clearpage
\null\vfill
\begin{center}
\phantomsection\captionof{figure}{The cross-subsidy under a break-even flat fee, by state}\label{fig-cross-area}
\par\vspace{8pt}
\fbox{\textit{[pending figure: Output/Figures/Fig\_CrossSub\_Area\_Panel]}}
\par\vspace{8pt}
\begin{minipage}{0.85\textwidth}\small\textit{Notes:} The break-even fee is the average fair premium, the single flat fee that exactly covers the fund's cleanups. Facilities below the line overpay (blue); facilities above it are subsidized (red); the two areas are equal. The transfer is 26 percent of premiums in New Mexico, 18 in Colorado, 18 in Louisiana, and 16 in Tennessee. 2005 cross-section.\end{minipage}
\end{center}
\vfill\clearpage

\clearpage
\null\vfill
\begin{center}
\phantomsection\captionof{figure}{The cross-subsidy over time, Tennessee}\label{fig-cross-time}
\par\vspace{8pt}
\fbox{\textit{[pending figure: Output/Figures/Fig\_CrossSub\_AreaTime\_TN]}}
\par\vspace{8pt}
\begin{minipage}{0.85\textwidth}\small\textit{Notes:} Tennessee in 2000, 2005, 2010, and 2015. The break-even fee rises as the fleet ages, from \$3,845 to \$5,068, while the cross-subsidy falls as a share of premiums, from 18 to 15 percent, because facilities grow more alike. Colorado, Louisiana, and New Mexico appear in the same form by state; New Mexico is too thin to read a trend.\end{minipage}
\end{center}
\vfill\clearpage

\subsection{---- Section 4.4: does the policy work
----}\label{section-4.4-does-the-policy-work--}

\clearpage
\null\vfill
\begin{center}
\phantomsection\captionof{table}{Difference-in-differences estimates of the tank-closure response}\label{tbl-stepped-did}
\vspace{6pt}
\fbox{\textit{[pending table: Output/Tables/T\_Stepped\_DiD\_OLS]}}
\par\vspace{10pt}
\begin{minipage}{0.85\textwidth}\small\textit{Notes:} Effect of the 1998 Texas reform on the annual tank-closure indicator. The preferred column adds facility and cell-by-year fixed effects. Standard errors clustered by state in parentheses.\end{minipage}
\end{center}
\vfill\clearpage

\clearpage
\null\vfill
\begin{center}
\phantomsection\captionof{figure}{Event study of tank closure, Texas relative to control states}\label{fig-es-tank}
\par\vspace{8pt}
\fbox{\textit{[pending figure: Output/Figures/Fig\_ES\_HTE\_Pooled]}}
\par\vspace{8pt}
\begin{minipage}{0.85\textwidth}\small\textit{Notes:} Year-by-year closure gap between Texas and control tanks, relative to the year before the reform. Pre-reform coefficients are flat; the gap opens at the reform and persists.\end{minipage}
\end{center}
\vfill\clearpage

\clearpage
\null\vfill
\begin{center}
\phantomsection\captionof{table}{Risk concentration: closure response by tank age and wall type}\label{tbl-risk-conc}
\vspace{6pt}
\fbox{\textit{[pending table: Output/Tables/T\_Falsification\_RiskConcentration]}}
\par\vspace{10pt}
\begin{minipage}{0.85\textwidth}\small\textit{Notes:} The closure response interacted with tank age and wall type, with facility and year fixed effects placed outside the make-model cell so age and wall are not absorbed. The response concentrates in old single-walled tanks.\end{minipage}
\end{center}
\vfill\clearpage

\clearpage
\null\vfill
\begin{center}
\phantomsection\captionof{table}{Facility-level portfolio response to the reform}\label{tbl-fac-att}
\vspace{6pt}
\fbox{\textit{[pending table: Output/Tables/T\_Facility\_Portfolio\_ATT\_Pub]}}
\par\vspace{10pt}
\begin{minipage}{0.85\textwidth}\small\textit{Notes:} Facility-year outcomes under facility and portfolio-mix-by-year fixed effects, with the composition-weighted baseline as a control. Standard errors clustered by state.\end{minipage}
\end{center}
\vfill\clearpage

\clearpage
\null\vfill
\begin{center}
\phantomsection\captionof{figure}{Event study of facility-level closure, Texas relative to control states}\label{fig-es-fac}
\par\vspace{8pt}
\fbox{\textit{[pending figure: Output/Figures/Fig\_ES\_Facility\_Portfolio]}}
\par\vspace{8pt}
\begin{minipage}{0.85\textwidth}\small\textit{Notes:} Facility-level closure gap by year under the portfolio-mix-by-year design, relative to the year before the reform. Flat before the reform, rising after.\end{minipage}
\end{center}
\vfill\clearpage

\clearpage
\null\vfill
\begin{center}
\phantomsection\captionof{table}{Effect of risk-based pricing on leak discovery}\label{tbl-leak-did}
\vspace{6pt}
\fbox{\textit{[pending table: Output/Tables/T\_LUST\_A\_Total]}}
\par\vspace{10pt}
\begin{minipage}{0.85\textwidth}\small\textit{Notes:} Effect of the reform on the annual confirmed-leak-discovery indicator. First column is the pre-reform control mean; second is the Texas-times-post coefficient. Facility and cell-by-year fixed effects, clustered by state.\end{minipage}
\end{center}
\vfill\clearpage

\clearpage
\null\vfill
\begin{center}
\phantomsection\captionof{table}{Leak discovery by detection channel}\label{tbl-leak-decomp}
\vspace{6pt}
\fbox{\textit{[pending table: Output/Tables/T\_LUST\_B\_Decomp]}}
\par\vspace{10pt}
\begin{minipage}{0.85\textwidth}\small\textit{Notes:} The leak-discovery effect split into the routine-monitoring channel and the closure-inspection channel. Discovery falls through both, so the reform lowers the true leak rate rather than only its detection.\end{minipage}
\end{center}
\vfill\clearpage

\subsection{---- Section 4.5: who, where, and what
----}\label{section-4.5-who-where-and-what--}

\clearpage
\null\vfill
\begin{center}
\phantomsection\captionof{figure}{Event study of facility downsizing, Texas relative to control states}\label{fig-es-downsize}
\par\vspace{8pt}
\fbox{\textit{[pending figure: Output/Figures/Fig\_ES\_Facility\_Downsize]}}
\par\vspace{8pt}
\begin{minipage}{0.85\textwidth}\small\textit{Notes:} Year-by-year gap in facility downsizing (fewer tanks and fewer gallons), relative to the year before the reform. The pre-reform path is not flat, so I read this margin as descriptive.\end{minipage}
\end{center}
\vfill\clearpage

\clearpage
\null\vfill
\begin{center}
\phantomsection\captionof{figure}{Event study of facility consolidation, Texas relative to control states}\label{fig-es-consolidate}
\par\vspace{8pt}
\fbox{\textit{[pending figure: Output/Figures/Fig\_ES\_Facility\_Consolidate]}}
\par\vspace{8pt}
\begin{minipage}{0.85\textwidth}\small\textit{Notes:} Year-by-year gap in facility consolidation (fewer tanks, about the same gallons), relative to the year before the reform. The path is flat and noisy with no clear break at the reform, consistent with the null average effect.\end{minipage}
\end{center}
\vfill\clearpage

\clearpage
\null\vfill
\begin{center}
\phantomsection\captionof{figure}{Event study of facility replacement, Texas relative to control states}\label{fig-es-replace}
\par\vspace{8pt}
\fbox{\textit{[pending figure: Output/Figures/Fig\_ES\_Facility\_Replace]}}
\par\vspace{8pt}
\begin{minipage}{0.85\textwidth}\small\textit{Notes:} Year-by-year gap in facility replacement (tanks swapped at about the same capacity), relative to the year before the reform. The pre-reform path is acceptable and the response is moderate.\end{minipage}
\end{center}
\vfill\clearpage

\clearpage
\null\vfill
\begin{center}
\phantomsection\captionof{table}{Closure response by fuel type (tank level)}\label{tbl-hte-tank}
\vspace{6pt}
\fbox{\textit{[pending table: Output/Tables/T\_HTE\_Tank\_Pub]}}
\par\vspace{10pt}
\begin{minipage}{0.85\textwidth}\small\textit{Notes:} Triple-difference interactions of the treatment with fuel type, with facility and cell-by-year fixed effects. The response concentrates in motor-fuel-only facilities.\end{minipage}
\end{center}
\vfill\clearpage

\clearpage
\null\vfill
\begin{center}
\phantomsection\captionof{table}{Facility response by fuel, place, and market (heterogeneity)}\label{tbl-fac-hte}
\vspace{6pt}
\fbox{\textit{[pending table: Output/Tables/T\_Facility\_HTE\_closure\_share\_Pub]}}
\par\vspace{10pt}
\begin{minipage}{0.85\textwidth}\small\textit{Notes:} Facility-level triple-difference interactions of the treatment with facility attributes fixed at the reform. Gas stations carry the response across margins.\end{minipage}
\end{center}
\vfill\clearpage

\clearpage
\null\vfill
\begin{center}
\phantomsection\captionof{table}{Portfolio response by facility size}\label{tbl-size-hte}
\vspace{6pt}
\fbox{\textit{[pending table: Output/Tables/T\_Facility\_SizeHTE\_Pub]}}
\par\vspace{10pt}
\begin{minipage}{0.85\textwidth}\small\textit{Notes:} Treatment interacted with total-capacity bin, reference the smallest bin. Small facilities exit; large facilities downsize and replace. Facility and portfolio-mix-by-year fixed effects, clustered by state.\end{minipage}
\end{center}
\vfill\clearpage

\clearpage
\null\vfill
\begin{center}
\phantomsection\captionof{table}{Portfolio response by tank vintage}\label{tbl-vintage-hte}
\vspace{6pt}
\fbox{\textit{[pending table: Output/Tables/T\_Facility\_VintageHTE\_Pub]}}
\par\vspace{10pt}
\begin{minipage}{0.85\textwidth}\small\textit{Notes:} Treatment interacted with installation cohort, reference the 1989 to 1998 cohort. The oldest tanks exit without reinvesting; the youngest reinvest. Facility and portfolio-mix-by-year fixed effects, clustered by state.\end{minipage}
\end{center}
\vfill\clearpage

\subsection{---- Section 6: structural estimates
----}\label{section-6-structural-estimates--}

\clearpage
\null\vfill
\begin{center}
\phantomsection\captionof{table}{Structural parameter estimates}\label{tbl-struct-estimates}
\vspace{6pt}
\fbox{\textit{[pending table: Output/Tables/T\_Struct\_Estimates]}}
\par\vspace{10pt}
\begin{minipage}{0.85\textwidth}\small\textit{Notes:} Nested pseudo-likelihood estimates of the dynamic portfolio model. Values are in model units of \$10,000 per facility-year. Headline pooled-exposure coefficients pending the converged run.\end{minipage}
\end{center}
\vfill\clearpage


  \bibliography{UST-lit-fixed.bib}


\end{document}
